\begin{table}[htb]
    \caption[Moderationsanalyse Religiosität und kulturelle Orientierung auf das psychische Wohlbefinden]{\textit {Moderationsanalyse Religiosität und kulturelle Orientierung auf das psychische Wohlbefinden}} 
    \label{Moderationsanalyse}
    \centering
    \begin{adjustbox}{width=11cm} %{width=\textwidth}
    \small
    \begin{tabular}{lrrrrrr}
      \hline
                 & $B$   & $SE(B)$&  $t$    & $p$   & $\Delta R^{2}$\\
      \hline
    Konstante    & 2.24    & 0.45 & 4.95    & $<$~.001 & \\
    Religiosität & $-$0.47 & 0.41 & $-$1.15 & .250     & \\
    Kultur       & 0.21    & 0.10 & 2.16    & .032     & \\
    Interaktionsterm & 0.14 & 0.09 & 1.56   & .119     & .006 \\
    $R^{2}$          &      &      &        &          & .063* \\
       \hline
    \end{tabular}
    \end{adjustbox}
    
    \begin{tablenotes}
        \item \textit{Anmerkungen.} \(N\)~=~418; Wertebereich der Variable Religiosität -1.69 (\textit{gar nicht religiös}) bis 1.16 (\textit{sehr religiös}); kulturelle Orientierung 1 (\textit{überhaupt nicht}) bis 7 (\textit{immer}).\\ *$p<$~.001
      \end{tablenotes}
    \end{table}


% Wertebereich Religiosität muss noch
